% =====================================================================
% Capturing Provenance of Academy Examples with Flowcept -- 1h tutorial
% Compile: make   (pdflatex + beamer + metropolis)
% Outline: Academy reminder -> Flowcept + ex 1/2/3 -> agents & LLMs
%          -> advanced: LangGraph (07) + Parsl (05) -> inspecting failures
% =====================================================================
\documentclass[aspectratio=169,11pt]{beamer}
\usetheme{metropolis}
\usepackage{listings}
\usepackage{xcolor}
\usepackage{booktabs}

\definecolor{fcblue}{HTML}{2A78D6}
\definecolor{fcorange}{HTML}{EB6834}
\definecolor{fcaqua}{HTML}{1BAF7A}
\definecolor{fcink}{HTML}{0B0B0B}
\definecolor{fcmuted}{HTML}{898781}
\definecolor{fccode}{HTML}{F5F5F2}
\definecolor{fcred}{HTML}{D03B3B}

\setbeamercolor{frametitle}{bg=fcblue}
\setbeamercolor{progress bar}{fg=fcorange}
\setbeamercolor{alerted text}{fg=fcorange}

\lstdefinestyle{py}{language=Python,basicstyle=\ttfamily\scriptsize,
  keywordstyle=\color{fcblue}\bfseries,commentstyle=\color{fcmuted}\itshape,
  stringstyle=\color{fcaqua},backgroundcolor=\color{fccode},
  numbers=none,frame=none,breaklines=true,showstringspaces=false,
  aboveskip=3pt,belowskip=2pt,morekeywords={async,await,self,with,as,for}}
\lstset{style=py}
\lstdefinestyle{sh}{language=bash,basicstyle=\ttfamily\scriptsize,
  backgroundcolor=\color{fccode},breaklines=true}

\title{Capturing Provenance of Academy Agents with Flowcept}
\subtitle{Hands-on tutorial for ATPESC 2024}
\author{Amal Gueroudji}
\date{ATPESC}
\institute{Argonne National Laboratory}

\begin{document}

\maketitle

% =====================================================================
\section{Provenance --- another kind of data}

\begin{frame}{Provenance: another kind of data}
  Alongside the inputs and results a workflow produces, there is \textbf{another kind
  of data} worth keeping: \alert{provenance} --- the recorded history of \emph{what
  ran, on what, producing what, when, where, and how it all connects}.
  \vspace{0.5em}
  \begin{itemize}
    \item It is newly \textbf{essential} because of AI: \alert{agentic workflows are
      non-deterministic} --- the same prompt yields different outputs, and the control
      flow is \emph{decided at runtime}, not written in your source.
    \item Without that record an AI run is \emph{not reproducible} and its failures are
      \emph{invisible} --- you cannot say what the agent actually did.
  \end{itemize}
  \vspace{0.3em}
  \centering In this session we \textbf{capture and study provenance with
  \alert{Flowcept}}, on agentic workflows built with \alert{Academy}.
\end{frame}

\begin{frame}{Why capture provenance at all?}
  Provenance = the recorded history of \emph{what ran, on what inputs, producing
  what outputs, when, where, and how it all connects}.
  \vspace{0.5em}
  \begin{itemize}
    \item \textbf{Reproducibility} --- rerun and get the same result; know exactly
      what code, data, and parameters produced a number.
    \item \textbf{Debugging} --- when something breaks, see \emph{which} step failed,
      on \emph{what} input, and why --- not just a stack trace.
    \item \textbf{Trust \& audit} --- verify and defend results; a scientific claim
      is only as good as the record behind it.
    \item \textbf{Cost \& performance} --- see where time and tokens go.
  \end{itemize}
  \vspace{0.3em}
  \small For classic HPC workflows this is valuable. For \alert{agentic} systems it
  becomes \alert{essential} --- next slide.
\end{frame}

\begin{frame}{Why it matters \emph{more} for agentic systems}
  Traditional code does what you wrote. \textbf{Agents decide what to do at
  runtime} --- so the ``what happened'' isn't in your source.
  \begin{itemize}\footnotesize
    \item \textbf{Non-determinism.} LLM calls sample --- you can't reproduce a run
      without the actual prompt, response, model, and tokens.
    \item \textbf{Autonomy \& emergence.} \texttt{@loop}s act on their own; agents
      choose which tool/peer to call. The control flow is \emph{decided}, not
      declared --- capture is the only record of the real path.
    \item \textbf{Distribution.} Work spans agents, processes, nodes, and
      \emph{frameworks} (Academy + LangGraph + Parsl); lineage across those
      boundaries is invisible without provenance.
    \item \textbf{Opaque reasoning.} Prompts, tool calls, and delegations bury the
      ``why.'' Provenance keeps every prompt/response.
    \item \textbf{Failure is normal.} Small LLMs propose malformed inputs; agents
      derail mid-plan --- you need the exact failing step + input.
  \end{itemize}
  \centering\footnotesize\alert{Non-deterministic, autonomous, distributed ---
  exactly the case provenance was made for.}
\end{frame}

% =====================================================================
\section{The tools --- Flowcept + Academy}

\begin{frame}{What is Flowcept?}
  \textbf{Flowcept} is a runtime \textbf{provenance capture} system: it records
  \emph{what ran, with what inputs/outputs, when, and how it connects} --- the
  provenance graph.
  \vspace{0.3em}
  \begin{itemize}
    \item \textbf{Data model}: \texttt{campaign} $\to$ \texttt{workflow} $\to$
      \texttt{task} (action / loop / llm\_call), with \texttt{used}/\texttt{generated}
      and \texttt{parent\_task\_id} lineage.
    \item Its \emph{agentic} plugin patches Academy's \texttt{Runtime}, so capture is
      \alert{automatic --- zero changes to agent logic}.
    \item Runs \textbf{offline} here (no Redis / no Mongo): records land in a per-run
      \texttt{flowcept\_buffer.jsonl}.
  \end{itemize}
\end{frame}

\begin{frame}{What is Academy?}
  A framework for \textbf{distributed multi-agent systems} for science --- these are
  the agentic workflows we instrument (unchanged).
  \begin{itemize}
    \item \texttt{Agent} --- a class holding state + behaviors
    \item \texttt{@action} --- a callable behavior (request/response)
    \item \texttt{@loop} --- autonomous background behavior
    \item \texttt{Handle} --- a remote reference to an agent
    \item \texttt{Manager} + \texttt{Exchange} --- launch agents \& route messages
  \end{itemize}
  \vspace{0.3em}
  \small Scales from local threads $\to$ processes $\to$ Globus Compute / Parsl across nodes.
\end{frame}

\begin{frame}[fragile]{Academy: the shape of an agent}
\begin{lstlisting}
class Coordinator(Agent):
    @action
    async def process(self, text: str) -> str:
        text = await self.lowerer.lower(text)      # cross-agent call
        return await self.reverser.reverse(text)   # cross-agent call

async with await Manager.from_exchange_factory(
        factory=LocalExchangeFactory(),
        executors=ThreadPoolExecutor()) as manager:
    coord = await manager.launch(Coordinator, args=(lo, re))
    await coord.process("DEADBEEF")     # -> "feebdaed"
\end{lstlisting}
  \small These are Academy's own examples
  (\texttt{github.com/academy-agents/academy}); we instrument them unchanged.
\end{frame}

% =====================================================================
\section{The hands-on session}

\begin{frame}{The hands-on session}
  Take the \textbf{upstream Academy examples}, instrument each with \textbf{Flowcept}
  to capture provenance (\alert{zero changes to agent logic}), and \textbf{inspect}
  what happened --- including \alert{where things fail}.
  \vspace{0.6em}
  \begin{tabular}{r l}
    \toprule
    \textbf{Part} & \textbf{Content} \\
    \midrule
    1 & Provenance --- another kind of data, and why agentic makes it essential \\
    2 & The tools: Flowcept + Academy \\
    3 & Flowcept + examples 1, 2, 3 --- and how to inspect provenance \\
    4 & Advanced: agents \& LLMs (6, 8), LangGraph + real xTB (7), Parsl (5), failures \\
    \bottomrule
  \end{tabular}
  \vspace{0.6em}\\
  \small \textbf{Terminal-only} (no images): text summaries, ASCII dashboards, a
  markdown provenance card --- runs over SSH on Aurora. Every run writes to its own
  \texttt{runs/<id>\_<date-time>/}.
\end{frame}

\begin{frame}[fragile]{How the tutorial is laid out}
  Eight upstream Academy examples, each a \textbf{step-by-step exercise}. As shipped
  each one \emph{runs but records nothing}; you turn provenance on one \textbf{STEP}
  at a time by uncommenting a block and re-running.
\begin{lstlisting}[style=sh]
exercises/
  local/    run on a laptop / login shell   (python exercise.py)
  aurora/   run on ALCF Aurora              (qsub submit.pbs; shared env.sh)
    <id>/   exercise.py  solution.py  agent_src.py  README.md  [submit.pbs]
\end{lstlisting}
  \begin{itemize}\small
    \item \texttt{exercise.py} --- STEP 0 baseline, then STEP 1 capture, 2 inspect,
      3 analyze, 4 card, 5 query (uncomment as you go)
    \item \texttt{solution.py} --- every step already enabled (the reference)
    \item \texttt{flowcept\_academy/} is a small \textbf{library} the exercises build on
  \end{itemize}
\end{frame}

\begin{frame}[fragile]{Setup --- one command, one env, from scratch}
\begin{lstlisting}[style=sh]
bash setup/install.sh        # builds ONE conda env for all 8 (CPU)
conda activate flowcept-academy
cd exercises/local/01-actor-client && python exercise.py
\end{lstlisting}
  \small \texttt{install.sh} builds a single \texttt{flowcept-academy} conda env and
  installs Flowcept (agentic branch) + Academy + Parsl + LangGraph + a CPU build of
  \texttt{torch}/\texttt{transformers}, plus \texttt{rdkit}/\texttt{xtb-python} for
  example 07's real chemistry, and drops the offline Flowcept settings in place.
  \alert{One env for the whole tutorial} --- no venv, no per-exercise env. No Redis,
  no Mongo, no GPU.
\end{frame}

% =====================================================================
\section{Capture + inspect --- examples 1, 2, 3}

\begin{frame}[fragile]{Turn capture on --- STEP 1}
\begin{lstlisting}
from flowcept_academy.capture import captured   # zero agent-code changes
from flowcept_academy.util import capture_run    # -> runs/<id>_<date-time>/

with capture_run("03-agent-agent") as run_dir:
    with captured(workflow_name="03-agent-agent"):
        result = asyncio.run(run())    # the upstream example, now captured
    # -> flowcept_buffer.jsonl in run_dir  (offline: no Redis / no Mongo)
\end{lstlisting}
  \small In the exercise you just uncomment this block and re-run. Everything a run
  produces --- buffer, perf CSV, the markdown card --- lands together in the run dir:
\begin{lstlisting}[style=sh]
cd exercises/local/03-agent-agent && python exercise.py
# -> runs/03-agent-agent_<date-time>/ : flowcept_buffer.jsonl + card.md + perf.csv
\end{lstlisting}
\end{frame}

\begin{frame}{Outcomes: examples 1, 2, 3}
  Each example exercises a different slice of the provenance model:
  \vspace{0.4em}
  \begin{tabular}{l l}
    \toprule
    \textbf{Example} & \textbf{Provenance captured} \\
    \midrule
    01 actor-client & agent lifecycle + actions (\texttt{increment} $\to$ count) \\
    02 agent-loop & autonomous \texttt{@loop} start/exit events (\texttt{group\_id}) \\
    03 agent-agent & \textbf{cross-agent calls} (caller $\to$ callee.action) \\
    \bottomrule
  \end{tabular}
  \vspace{0.6em}\\
  \small e.g. 03 shows the chain \texttt{DEADBEEF} $\to$ \texttt{deadbeef} $\to$
  \texttt{feebdaed} across three agents.
\end{frame}

\begin{frame}[fragile]{How to inspect the provenance}
  Three built-in ways, for \emph{any} run --- all terminal (from an example folder):
\begin{lstlisting}[style=sh]
# 1) tailored, content-aware analysis + text dashboard (printed by each run)
python solution.py                       # STEPs 2-4 print the report

# 2) analyze any captured buffer / whole run dir
python ../../../provenance/analyze.py runs/03-agent-agent_*

# 3) interactive queries: ask("...") turns NL into pandas
python ../../../provenance/query.py runs/03-agent-agent_* --ask "show the cross-agent calls"
\end{lstlisting}
  \small The analysis auto-adapts: it reports only what the example produced
  (actions, loops, cross-agent calls, LLM calls, \dots) --- and any failures.
  Flowcept's markdown provenance card is written into each run dir.
\end{frame}

\begin{frame}[fragile]{The provenance dashboard --- in the terminal}
  No images: each run prints an ASCII dashboard (records by subtype; LLM tokens by
  agent, or tasks by activity; capture overhead). \small Real \texttt{06-llm} run:
\begin{lstlisting}[style=sh]
============================================================
06-llm -- provenance
============================================================
  records by subtype:
    academy_lifecycle  | ################################ 4
    llm_call           | ################                 2
    academy_action     | ################                 2
  LLM tokens by agent:
    Orchestrator       | ################################ 140
  capture overhead (total ms by category):
    flush              | ################################ 1.135
    llm_hook           | #########################        0.892
    intercept_task     | #######################          0.836
============================================================
\end{lstlisting}
\end{frame}

% =====================================================================
\section{Advanced --- agents \& LLMs}

\begin{frame}{LLM backend: Argo $\to$ vLLM $\to$ OpenAI $\to$ local CPU model}
  \begin{itemize}\small
    \item \texttt{ARGO\_USER} set $\to$ \alert{ANL Argo gateway} (on Aurora / ANL network), native tool calling
    \item else \texttt{VLLM\_BASE\_URL}/\texttt{OPENAI\_BASE\_URL} set $\to$ \alert{vLLM}
      server we run (\texttt{VLLM\_MODEL}); native tool calling when launched with
      \texttt{--enable-auto-tool-choice} (on Aurora: \texttt{source ../vllm\_serve.sh \&\& vllm\_start})
    \item else \texttt{OPENAI\_API\_KEY} set $\to$ \alert{OpenAI}
      (\texttt{api.openai.com}, \texttt{OPENAI\_MODEL}, default \texttt{gpt-4o-mini}), native tool calling
    \item else $\to$ \alert{local CPU model}
      (\texttt{Qwen2.5-0.5B-Instruct}, offline, no GPU)
    \item \textbf{no mock}: every response is from a real model; \texttt{chat()}
      raises rather than fabricate text
  \end{itemize}
  \small Checked in priority order, first match wins; force with
  \texttt{FLOWCEPT\_TUTORIAL\_LLM=argo|vllm|openai|local}. Every call is recorded as an
  \texttt{llm\_call} (model, prompt, response, tokens), linked to its action.
\end{frame}

\begin{frame}[fragile]{Examples 6 \& 8: LLM-driven agents (vendored upstream)}
\begin{lstlisting}
# 06: the upstream Orchestrator builds a LangChain ReACT agent over a
# tool that messages a separate MySimAgent -- the LLM decides to call it.
agent = create_agent(make_chat_model(), tools=[sim_tool])
result = await agent.ainvoke({"messages": [("user", question)]})
\end{lstlisting}
  \begin{itemize}\small
    \item \textbf{06 llm}: ReACT agent $\to$ \texttt{llm\_call} +
      \texttt{langgraph\_node}/\texttt{graph}, and a cross-agent tool call when the
      LLM invokes the tool
    \item \textbf{08 discussion}: multi-agent group chat --- three
      \texttt{GroupChatAgent}s (Manager, Assistant, Senior Engineer) +
      \texttt{RoundRobinGroupChatManager} with a supervisor \texttt{@loop};
      per-turn \texttt{llm\_call}s, tokens per agent, bounded by \texttt{max\_rounds}
  \end{itemize}
  \small The LLM is injected via \texttt{make\_chat\_model()} (Argo $\to$ vLLM $\to$ OpenAI
  $\to$ local) exactly where upstream builds \texttt{ChatOpenAI} --- \alert{zero other agent-code changes}.
\end{frame}

\begin{frame}[fragile]{What gets captured --- a real \texttt{06-llm} run}
  \small The upstream ReACT agent, run on the \textbf{local 0.5B} model. One
  \texttt{@action} wraps one \texttt{langgraph\_graph}/\texttt{node}, which drives one
  \texttt{llm\_call} --- all sharing the campaign:
\begin{lstlisting}[style=sh]
* Records (subtype  activity)
    academy_action    answer               <- Orchestrator.@action
    langgraph_graph   LangGraph            <- create_agent ReACT graph
    langgraph_node    model                <- the LLM step in the graph
    llm_call          _AsyncChatHuggingFace
        prompt : ...simulated ionization energy of benzene...
        answer : The simulated ionization energy of benzene can be
                 calculated using quantum chemistry methods, such as DFT...
\end{lstlisting}
  \small The action, the graph, its node, and the LLM call are one connected,
  queryable lineage. \alert{Note}: the 0.5B model answered \emph{directly} instead of
  calling the tool --- provenance is exactly how you \emph{see} that it skipped the
  tool. A tool-capable backend (Argo) instead emits a cross-agent
  \texttt{tool\_call} to \texttt{MySimAgent}.
\end{frame}

% =====================================================================
\section{Extra advanced --- failures}

\begin{frame}[fragile]{Example 7: cross-framework (Academy + LangGraph + real xTB)}
  The upstream \texttt{mol\_design\_agents.py}, \alert{unchanged}: a LangGraph
  \texttt{StateGraph} (plan $\to$ tool\_calling $\to$ simulate $\to$ conclude $\to$
  critique $\to$ update, looping) runs inside an Academy \texttt{@loop}, and
  \texttt{simulate} calls \textbf{real GFN2-xTB} (RDKit $\to$ ASE $\to$ xtb).
  \begin{itemize}\small
    \item \textbf{one shared \texttt{campaign\_id}} across both plugins (Academy + LangGraph)
    \item three layers in one lineage: \texttt{academy\_loop} $\to$
      \texttt{langgraph\_node} $\to$ \texttt{tool\_call} (the xTB energy)
    \item LLM = Argo \texttt{gpt-4o} (tool-capable) --- drives the \texttt{tool\_calling} node
  \end{itemize}
  \small Real run --- \texttt{tool\_call}=10, \texttt{langgraph\_node}=9,
  \texttt{llm\_call}=6, \texttt{academy\_loop}=2 (1 campaign):
\begin{lstlisting}[style=sh]
* nodes: plan, tool_calling(2), simulate(2), should_continue,
         conclude, critique, update
* real xTB ionization energies (charged - neutral, eV):
    simulate -> 13.1975...    simulate -> 14.1156...
\end{lstlisting}
\end{frame}

\begin{frame}[fragile]{Example 5: Parsl}
\begin{lstlisting}
class SimulationAgent(Agent):
    @action
    async def run_expensive_task(self) -> int:
        return await asyncio.wrap_future(expensive_task())  # a Parsl app
\end{lstlisting}
  \small The agent delegates to a \textbf{Parsl} task; provenance captures the
  dispatching action and its result. (Example 07 shows the same pattern with real
  chemistry --- the agent's process pool runs GFN2-xTB, each emitting a
  \texttt{tool\_call}.)
\end{frame}

\begin{frame}[fragile]{Inspecting failures --- where provenance shines}
  Agents (and LLMs) get things wrong. In \textbf{example 07 this is not staged} ---
  the LLM proposes candidate molecules and some simply \alert{can't be built} by
  RDKit/xTB. Flowcept records each failed \texttt{tool\_call} with
  \textcolor{fcred}{\texttt{status=ERROR}} + real \texttt{stderr}:
\begin{lstlisting}[style=sh]
  status:  FINISHED 22   ERROR 10
  !  Failures captured (status=ERROR) -- what the agent actually tried:
    tool_call/compute_ionization_energy  [ERROR]
      stderr: Parse failure for C3N2O
    tool_call/compute_ionization_energy  [ERROR]
      stderr: Parse failure for CF3NC(N)=O
    tool_call/compute_ionization_energy  [ERROR]
      stderr: Parse failure for C6H5NC(N)=O
    ... (8 xTB parse failures + 1 propagated langgraph_node)
\end{lstlisting}
  \small You see \emph{which} molecule the LLM invented, that xTB rejected it, and that
  successful runs still gave real energies (13.20, 14.12 eV) --- the debugging payoff.
\end{frame}

\begin{frame}[fragile]{Running on Aurora (offline, terminal-only)}
  Each \texttt{exercises/aurora/<id>/} has its own \texttt{submit.pbs}
  (\texttt{-A ATPESC2026}); all share \texttt{env.sh} (modules + conda + offline LLM).
\begin{lstlisting}[style=sh]
export FLOWCEPT_ENV_PREFIX=/lus/flare/projects/ATPESC2026/prov/envs/flowcept-academy
bash setup/install.sh aurora     # frameworks base + delta, no LLM download
source exercises/aurora/env.sh   # modules + conda --stack + HF offline
cd exercises/aurora/06-llm && qsub submit.pbs   # submit.pbs starts vLLM on the node
\end{lstlisting}
  \small \texttt{install.sh aurora} adds a small delta on the \texttt{frameworks} module
  (nothing downloaded); all LLM usage reads ALCF's staged hub offline
  (\texttt{HF\_HOME}=\texttt{/flare/datasets/model-weights}, \texttt{HF\_HUB\_OFFLINE=1}).
  \textbf{vLLM is the only backend} --- each \texttt{submit.pbs} starts it on the node
  (\alert{no CPU fallback}; \texttt{FLOWCEPT\_USE\_VLLM=0} $\to$ Argo). Details next slide.
\end{frame}

\begin{frame}[fragile]{vLLM on Aurora: serve it yourself (+ interactive)}
  \small On Aurora's XPU, \texttt{vllm serve} normally \alert{SIGSEGVs} inspecting the
  model architecture. \texttt{vllm\_start} dodges it: it \textbf{primes the modelinfo
  cache in-process} (model \emph{class} only --- no weights, no subprocess), then serves
  against the cache. \alert{Automatic}, nothing to do by hand.
  \begin{itemize}\footnotesize
    \item \textbf{Adopt-on-source}: sourcing \texttt{vllm\_serve.sh} in a fresh shell
      finds a running server + re-exports \texttt{VLLM\_BASE\_URL} (so \texttt{query.py
      --ask} routes to it).
    \item \textbf{Context}: \texttt{VLLM\_MAX\_MODEL\_LEN=32768} (8192 overflowed ---
      08's group chat accumulates the whole discussion).
  \end{itemize}
\begin{lstlisting}[style=sh]
qsub -I -A ATPESC2026 -q debug -l select=1 -l walltime=60:00   # interactive node
source ../env.sh                            # modules + conda + offline LLM
source ../vllm_serve.sh && vllm_start       # primes cache, serves on the GPUs
python solution.py                          # or: query.py runs/<id>_* --ask "..."
vllm_stop                                   # tear the server down; then exit
\end{lstlisting}
\end{frame}

\begin{frame}{Recap}
  \begin{itemize}
    \item \textbf{Capture:} one line (\texttt{captured()}) instruments any Academy
      example --- zero agent-code changes.
    \item \textbf{Examples:} actor, loop, cross-agent, multi-process, Parsl, LLM
      agents, and cross-framework LangGraph + real xTB (07).
    \item \textbf{LLM:} Argo $\to$ vLLM $\to$ OpenAI $\to$ local CPU model.
    \item \textbf{Inspect:} tailored analysis, text dashboards, provenance cards
      --- all terminal, including \alert{failures} (status=ERROR + stderr).
    \item \textbf{Outputs:} every run in its own \texttt{runs/<id>\_<date-time>/}.
  \end{itemize}
  \vspace{0.4em}
  \centering\Large\textbf{instrument $\to$ capture $\to$ inspect (even the failures)}
\end{frame}

\begin{frame}[standout]
  Thank you!\\[0.4em]
  \normalsize \texttt{exercises/\{local,aurora\}} · \texttt{provenance/analyze.py} ·
  \texttt{slides/} · upstream examples: \texttt{github.com/academy-agents/academy}
\end{frame}

\end{document}
